\documentclass[10pt]{amsart}

\usepackage[T1]{fontenc}
\usepackage{lmodern}
\usepackage[protrusion=true,expansion=false]{microtype}
\usepackage{mathtools}
\usepackage{amssymb}
\usepackage{mathrsfs}
\usepackage{booktabs}
\usepackage{enumitem}
\usepackage{xcolor}
\usepackage[colorlinks=true,linkcolor=blue!45!black,citecolor=blue!45!black,urlcolor=blue!55!black]{hyperref}
\usepackage[nameinlink,capitalize,noabbrev]{cleveref}

\numberwithin{equation}{section}

\newtheorem{theorem}{Theorem}[section]
\newtheorem{proposition}[theorem]{Proposition}
\newtheorem{lemma}[theorem]{Lemma}
\newtheorem{corollary}[theorem]{Corollary}
\theoremstyle{definition}
\newtheorem{definition}[theorem]{Definition}
\theoremstyle{remark}
\newtheorem{remark}[theorem]{Remark}

\newcommand{\CC}{\mathbf C}
\newcommand{\QQ}{\mathbf Q}
\newcommand{\PP}{\mathbf P}
\newcommand{\GGm}{\mathbf G_m}
\newcommand{\IC}{\operatorname{IC}}
\newcommand{\Nm}{\operatorname{Nm}}
\newcommand{\Alt}{\operatorname{Alt}}
\newcommand{\Sym}{\operatorname{Sym}}
\newcommand{\Stab}{\operatorname{Stab}}
\newcommand{\Sing}{\operatorname{Sing}}
\newcommand{\id}{\operatorname{id}}
\newcommand{\reg}{\mathrm{reg}}
\newcommand{\exc}{\mathrm{exc}}
\newcommand{\cH}{\mathcal H}
\newcommand{\cI}{\mathcal I}
\newcommand{\cM}{\mathcal M}
\newcommand{\cN}{\mathcal N}
\newcommand{\cB}{\mathcal B}
\newcommand{\cP}{\mathcal P}
\newcommand{\cX}{\mathcal X}
\newcommand{\cR}{\mathcal R}
\newcommand{\cU}{\mathcal U}
\newcommand{\cA}{\mathcal A}
\newcommand{\cC}{\mathcal C}
\newcommand{\cE}{\mathcal E}
\newcommand{\cL}{\mathcal L}
\newcommand{\wh}{\widehat}
\newcommand{\wt}{\widetilde}
\newcommand{\bo}{\boxtimes}

\title[A realization of $G_2$ for Litt Problem No. 13]
      {A Realization of $G_2$ for Litt Problem No. 13}
\author{OpenAI}
\date{Revised August 10, 2026}

\hypersetup{
  pdftitle={A Realization of G2 for Litt Problem No. 13},
  pdfauthor={OpenAI}
}

\subjclass[2020]{14K12, 14H40, 18M25, 20G41}
\keywords{abelian variety, cyclic cover, Prym variety, intersection complex,
convolution, Tannaka group, exceptional group $G_2$}

\begin{document}
\raggedbottom

\begin{abstract}
A family of complex abelian threefolds is constructed from cyclic triple covers
\[
  C:\quad w^3=y-c,\qquad c\ne0,
\]
of elliptic curves $y^2=x^3+Ax+B$.  The induced Prym polarization has type
$(1,1,3)$.  The centred theta section $X$ is a normal symmetric surface with
a unique simple elliptic singularity of type $\wt E_7$, and
$\chi(\IC_X)=7$.  For general parameters, the pair-incidence curve
$X\cap(-x-X)$ has exactly two nonexceptional components.  Their top
cohomology gives an unshifted summand $\IC_X\subset\Alt^2(\IC_X)$ in the
convolution category, while a generic unitary twist has Hodge multiplicities
$(2,3,2)$.  Together these facts determine the Tannaka pair: the connected
group is $G_2$ on its standard seven-dimensional representation, and the
same exterior-square morphism excludes the only possible disconnected
extension.  The full convolution Tannaka group is therefore $G_2$, realizing
the $G_2$ case of Litt Problem No.~13.
\end{abstract}

\maketitle

\section{Introduction}

Let $A$ be a complex abelian variety.  After negligible objects are
quotiented out, convolution turns the relevant constructible complexes on
$A$ into a neutral Tannakian category.  Thus an integral subvariety
$Z\subset A$ with $\chi(\IC_Z)\ne0$ determines a Tannaka group
$G(\IC_Z)$ and a faithful module of dimension $\chi(\IC_Z)$; see
\cite[Proposition~3.1 and Definition~3.2]{JKLM2025}.  Litt Problem No.~13 asks whether the four
exceptional simple types
\[
  G_2,\qquad E_7,\qquad E_8,\qquad F_4
\]
can occur in this manner \cite{Litt13}.  Only the $G_2$ case is considered
here; no claim is made about the other three.  The construction necessarily
uses a singular theta surface, so the exclusion results for smooth
subvarieties in \cite{KLM2026} do not apply.

The construction is defined over a rational parameter space.  Put
\[
 \Delta_E=4A^3+27B^2,
 \qquad
 \Delta_{\mathrm{br}}=4A^3+27(B-c^2)^2,
\]
and let
\[
 \cB_{\QQ}=\operatorname{Spec}
 \QQ[A,B,c,c^{-1},\Delta_E^{-1},\Delta_{\mathrm{br}}^{-1}],
 \qquad
 \cB=\cB_{\QQ}\times_{\QQ}\CC.
\]
For $b=(A,B,c)\in\cB(\CC)$, let
\[
 E_b:\ y^2=x^3+Ax+B,
 \qquad
 C_b:\ w^3=y-c.
\]
Here $C_b$ denotes the smooth projective cyclic cover defined by the section
$y-c\in H^0(E_b,\mathcal O_{E_b}(3O))$; the displayed equation is its
affine chart.  The reduced branch divisor makes this cover connected.  Let
\[
 \pi_b:C_b\longrightarrow E_b,
 \qquad (x,y,w)\longmapsto(x,y)
\]
be the cyclic triple cover.  The condition $\Delta_{\mathrm{br}}\ne0$
means that the line $y=c$ meets $E_b$ in three distinct points; these are
precisely the branch points of $\pi_b$.  Write $O\in E_b$ for the point at
infinity and set
\[
 \kappa_b=\pi_b^*\mathcal O_{E_b}(O),
 \qquad
 P_b=\ker(\Nm_{\pi_b})^0.
\]
The centred theta divisor on $J(C_b)$ restricts to the divisor
\[
 X_b=\{M\in P_b:h^0(C_b,\kappa_b\otimes M)>0\}.
\]
Write $\cP\to\cB$ and $\cX\subset\cP$ for the resulting relative Prym and
relative theta section.

\begin{theorem}[Main theorem]\label{thm:main}
There is a nonempty Zariski-open subset $\cB^\circ\subset\cB$ such that for
every $b\in\cB^\circ(\CC)$ the following hold.
\begin{enumerate}[label=\textup{(\roman*)},leftmargin=*]
\item $P_b$ is an abelian threefold, and the polarization induced by the
      centred theta divisor has type $(1,1,3)$.
\item $X_b\subset P_b$ is an integral normal symmetric ample Cartier divisor
      with a unique singular point at the origin; this singularity is simple
      elliptic of type $\wt E_7$.
\item $\chi(\IC_{X_b})=7$, and the associated Tannaka pair satisfies
\[
  \bigl(G(\IC_{X_b}),\omega(\IC_{X_b})\bigr)\simeq(G_2,V_7),
\]
      where $V_7$ is the standard seven-dimensional representation of
      $G_2$.  In particular, the full convolution Tannaka group is connected
      and isomorphic to $G_2$.
\end{enumerate}
Moreover, $\cX\subset\cP$ is a symmetric relatively ample effective Cartier
divisor, flat over $\cB$, with geometrically integral fibres.
\end{theorem}

\paragraph{Proof outline.}
Fix $b\in\cB(\CC)$ and write $(P,X)=(P_b,X_b)$.  The proof has a geometric
part and a Tannakian part, with the pair-incidence calculation linking the
two.  The norm fibre in $C^{(3)}$ is birational to $X$ and has one
$\wt E_6$ singularity.  Blowing up that point and contracting the strict
transform of the trigonal pencil gives the minimal resolution
$\rho:S\to X$; its exceptional locus is an elliptic curve of
self-intersection $-2$.  Adjunction, the local geometric genus, and Noether's
formula give
\[
 K_S^2=16,\qquad \chi(\mathcal O_S)=2,\qquad e(S)=8,
\]
while the theta group yields $p_g(S)=5$ and $q(S)=4$.  The semismall
decomposition then gives $\chi(\IC_X)=7$, and the unique singular point
forces $\Stab_P(X)=\{0\}$.  A generic unitary twist of the same resolution
has Hodge multiplicities $(2,3,2)$.

